%% ===========================================================================
%%  Propuesta de tesis — doble titulación Geología + Ingeniería Matemática
%%  "¿Qué se puede recuperar de una erupción?"
%%
%%  Compilar con:  latexmk -pdf propuesta_tesis.tex
%%             o:  pdflatex propuesta_tesis.tex   (dos veces)
%% ===========================================================================
\documentclass[aspectratio=169,10pt,t]{beamer}

\usepackage[T1]{fontenc}
\usepackage[utf8]{inputenc}
\usepackage[spanish,es-noshorthands,es-nodecimaldot]{babel}
\usepackage{lmodern}
\usepackage{amsmath,amssymb}
\usepackage{booktabs}
\usepackage{microtype}
\usepackage{tikz}
\usetikzlibrary{arrows.meta,positioning,decorations.pathmorphing,shapes.geometric}
\usepackage{appendixnumberbeamer}

\usetheme[progressbar=frametitle,numbering=fraction,block=fill]{metropolis}

\definecolor{magma}{HTML}{C1440E}
\definecolor{roca}{HTML}{2E4057}
\definecolor{gas}{HTML}{0E7C7B}
\definecolor{alerta}{HTML}{8B1E3F}

\setbeamercolor{progress bar}{fg=magma}
\setbeamercolor{title separator}{fg=magma}
\setbeamercolor{frametitle}{bg=roca}
\setbeamercolor{alerted text}{fg=magma}
\setbeamercolor{block title}{bg=roca!12,fg=roca}

% babel-spanish redefine \% con un chequeo de \lastskip que rompe en modo
% matemático ("Incompatible glue units"); \pc lo evita.
\newcommand{\pc}{\ifmmode\text{\,\%}\else\,\%\fi}

\newcommand{\MER}{\mathrm{MER}}
\newcommand{\dd}{\,\mathrm{d}}
\newcommand{\ideageo}{\textcolor{magma}{$\bigstar$}}
\newcommand{\fig}[2][\linewidth]{%
  \includegraphics[width=#1]{../prototipo/figuras/#2}}
\newcommand{\figh}[2][0.62\textheight]{%
  \includegraphics[height=#1,keepaspectratio]{../prototipo/figuras/#2}}

% Cajita de "mensaje de la lámina": el mensaje que se quiere que quede.
\newcommand{\mensaje}[1]{%
  \vfill
  \begin{beamercolorbox}[wd=\linewidth,sep=0.4em,rounded=true]{block title}
  \footnotesize #1
  \end{beamercolorbox}}

\title{¿Qué se puede recuperar de una erupción?}
\subtitle{Identificabilidad y no-unicidad en modelos de ascenso de magma por un conducto volcánico}
\date{Propuesta de tesis}
\author{Doble titulación: Geología + Ingeniería Matemática}
\institute{Caso de estudio: Calbuco, abril de 2015}

\begin{document}

\maketitle

% ===========================================================================
\begin{frame}{La reunión en una lámina}

\begin{columns}[T]
\begin{column}{0.48\textwidth}
\textbf{Dónde estamos}\medskip

Tenemos un modelo directo del conducto: dados los parámetros, calculamos
$(P,\phi,u_m,u_g)$ en $z$, en $(z,t)$ y en $(r,z)$.

\bigskip
\textbf{Qué propongo}\medskip

Recorrer la flecha al revés. Dados los observables de una erupción,
\alert{¿qué se puede afirmar sobre el magma en profundidad y sobre la geometría
del conducto, y con qué certeza?}
\end{column}

\begin{column}{0.48\textwidth}
\textbf{La tesis NO es}\medskip

``hacer una inversión y reportar números''.

\bigskip
\textbf{La tesis SÍ es}\medskip

caracterizar \alert{qué es recuperable y qué no}, porque aquí la no-unicidad no
es un accidente numérico: es \alert{estructural}, y se puede demostrar.

\bigskip
Ya tengo tres resultados preliminares que lo muestran.
\end{column}
\end{columns}

\mensaje{Y el hilo que une todo con el trabajo ya hecho: \textbf{agregar
dimensiones —tiempo, radio— es precisamente lo que permite recuperar parámetros
que el modelo estacionario 1D no puede distinguir.}}
\end{frame}

% ===========================================================================
\begin{frame}{Plan}
\setbeamertemplate{section in toc}[sections numbered]
{\small\tableofcontents[hideallsubsections]}
\end{frame}

% ===========================================================================
\section{El punto de partida}
% ===========================================================================

\begin{frame}{El sistema físico}
\begin{columns}[c]
\begin{column}{0.42\textwidth}
\centering
\begin{tikzpicture}[scale=0.80,every node/.style={font=\tiny}]
  % encajante
  \fill[roca!12] (-1.5,-0.85) rectangle (1.5,6.0);
  % conducto de radio variable
  \fill[magma!70]
    plot[smooth,domain=0:6,samples=40] ({0.42+0.45*exp(-\x/1.8)},\x) --
    (0.24,6) -- (-0.24,6) --
    plot[smooth,domain=6:0,samples=40] ({-0.42-0.45*exp(-\x/1.8)},\x) -- cycle;
  % cámara
  \fill[magma!70] (0,-0.05) ellipse (1.05 and 0.55);
  \draw[roca,thick] (0,-0.05) ellipse (1.05 and 0.55);
  % nivel de fragmentacion
  \draw[gas,very thick,dashed] (-1.45,4.5) -- (1.45,4.5);
  \node[gas,anchor=south west] at (0.45,4.55) {$z_f$: fragmentación};
  % vent
  \draw[roca,thick] (-1.5,6) -- (-0.24,6);
  \draw[roca,thick] (0.24,6) -- (1.5,6);
  \foreach \a in {62,76,90,104,118}
    \draw[gas!80,-{Latex[length=1.4mm]}] (0,6.05) -- ({1.05*cos(\a)},{6.05+1.05*sin(\a)});
  \node[gas,anchor=south] at (0,7.2) {$\MER(t)$, altura de columna};
  % etiquetas
  \node[anchor=north] at (0,-0.72) {cámara: $P(H)=P_{lit}+\Delta P$};
  \node[alerta,anchor=west] at (0.80,2.7) {$R(z)$ incógnita};
  \draw[alerta,-{Latex[length=1.2mm]}] (0.78,2.7) -- (0.58,2.7);
  \draw[roca,-{Latex[length=1.6mm]}] (-1.25,0.7) -- (-1.25,5.8);
  \node[roca,anchor=east] at (-1.30,3.3) {$z$};
\end{tikzpicture}
\end{column}

\begin{column}{0.55\textwidth}
Flujo bifásico ascendente por un conducto esbelto ($L/R\sim400$):
\medskip
\begin{itemize}
  \item la descompresión exsuelve agua $\Rightarrow$ burbujas;
  \item la desgasificación \alert{endurece} el fundido varios órdenes de
        magnitud (viscosidad de Giordano);
  \item al superar un umbral de porosidad, el magma \alert{fragmenta} y pasa de
        líquido con burbujas a gas con piroclastos.
\end{itemize}
\bigskip
Lo que se \emph{mide} está todo arriba: tasa de descarga másica, altura de
columna, texturas y densidades de los piroclastos, deformación del edificio,
sismicidad.

\medskip
Lo que se \emph{quiere saber} está todo abajo.
\end{column}
\end{columns}
\end{frame}

% ---------------------------------------------------------------------------
\begin{frame}{Qué hay construido hoy en el repositorio}
\small
\begin{tabular}{@{}p{0.23\textwidth}p{0.42\textwidth}p{0.28\textwidth}@{}}
\toprule
\textbf{Módulo} & \textbf{Formulación} & \textbf{Cierre} \\
\midrule
\texttt{RIconduitex5\_5} &
DAE de 6 variables $(P,\phi,N_d,x,u_m,u_g)$: momentum de dos fluidos con
arrastre, coalescencia, permeabilidad, cristalización cinética &
tiro sobre $v_{inicial}$; salida estrangulada o $P=P_{atm}$ \\[0.5em]
\texttt{RIconduit1D\_transient} &
3 EDP de transporte $(\phi,N_d,x)$ + presión cuasi-estática; IMEX
(RK4 + relajación exacta) &
bisección sobre $q$ para $P_{vent}=P_{atm}$ \\[0.5em]
\texttt{RIconduit2D\_FD} &
perfil radial $u_m(r)$ por diferencias finitas en cada nivel $z$ (lubricación),
Thomas + Picard; marcha RK4 en $z$ &
tiro sobre $v_{inicial}$ con $P_{exit}\approx P_{atm}$ \\
\bottomrule
\end{tabular}

\mensaje{Ya tenemos \textbf{tiempo} y \textbf{dimensión radial}. El trabajo de
simulación está hecho. La pregunta abierta es qué hacemos con él.}
\end{frame}

% ---------------------------------------------------------------------------
\begin{frame}{Tres obstáculos técnicos, honestamente}
\begin{enumerate}
  \item<1-> \textbf{Los tres solvers están desacoplados.} El 2D no usa el
    transiente ni viceversa; comparten sólo \texttt{calbuco2015d.py} y las
    constitutivas. Un modelo inverso necesita \alert{un} operador directo, no
    tres.
  \item<2-> \textbf{El operador directo no es diferenciable.} Hay
    $\max(0,\cdot)$, cambios de régimen con paradas por cruce de $\phi$,
    criterios de estrangulamiento y bisecciones anidadas. Eso descarta los
    métodos basados en gradiente hasta que se regularice.
  \item<3-> \textbf{Hay inconsistencias entre versiones.} Por ejemplo
    $\phi_{crit}$ dinámico (número de capilaridad) en el estacionario frente a
    un valor fijo $0.75$ en el transiente; el término de permeabilidad
    \texttt{\_Fmg} está definido pero no acoplado en el transiente.
\end{enumerate}

\onslide<4->{
\mensaje{Ninguno es un defecto del trabajo previo: son exactamente las tareas
que convierten un \textbf{código de simulación} en un \textbf{código de
inversión}, y constituyen el primer capítulo natural de la tesis.}}
\end{frame}

% ===========================================================================
\section{El giro: del problema directo al inverso}
% ===========================================================================

\begin{frame}{La flecha, en los dos sentidos}
\centering
\begin{tikzpicture}[every node/.style={font=\small},node distance=1.2cm]
  \node[draw,rounded corners,fill=roca!10,minimum width=3.4cm,minimum height=1.5cm,align=center]
    (th) {\textbf{Parámetros} $\theta$\\[0.2em]
          \scriptsize $R(\cdot),\ \Delta P,\ c_0,\ T,\ \xi,\ \phi_{crit},\ k(\phi)$};
  \node[draw,rounded corners,fill=roca!10,minimum width=3.4cm,minimum height=1.5cm,align=center,right=5.2cm of th]
    (d) {\textbf{Observables} $d$\\[0.2em]
         \scriptsize $\MER(t)$, masa, altura,\\ \scriptsize texturas, deformación, gas};
  \draw[-{Latex[length=3mm]},line width=1.1pt,roca] (th) -- node[above,yshift=2pt]{$\mathcal{F}$: \textbf{directo} — \emph{hecho}} node[below,yshift=-2pt,font=\scriptsize]{simular} (d);
  \draw[-{Latex[length=3mm]},line width=1.1pt,magma] ([yshift=-1.05cm]d.south) -- node[below,yshift=-2pt]{$\mathcal{F}^{-1}$: \textbf{inverso} — \emph{la tesis}} ([yshift=-1.05cm]th.south);
  \draw[magma] (d.south) -- ([yshift=-1.05cm]d.south);
  \draw[magma] (th.south) -- ([yshift=-1.05cm]th.south);
\end{tikzpicture}

\vspace{1em}
\begin{columns}[T]
\begin{column}{0.9\textwidth}
Las tres preguntas clásicas sobre $\mathcal{F}$, con
$d^{obs}=\mathcal{F}(\theta^\dagger)+\varepsilon$:
\medskip
\begin{enumerate}
  \item \textbf{Unicidad / identificabilidad.} ¿Es $\mathcal{F}$ inyectivo? Si
    no lo es, caracterizar $\mathcal{F}^{-1}(d)$.
  \item \textbf{Estabilidad.} ¿Es $\mathcal{F}^{-1}$ continuo? ¿Con qué módulo?
  \item \textbf{Reconstrucción.} Algoritmo, regularización, incertidumbre.
\end{enumerate}
\end{column}
\end{columns}

\mensaje{En este problema la respuesta a (1) es \alert{\textbf{no}}, y se puede
\textbf{demostrar}. Ése es el núcleo de la propuesta.}
\end{frame}

% ---------------------------------------------------------------------------
\begin{frame}{Las ideas que ya traíamos, ordenadas}
\small
\begin{tabular}{@{}p{0.40\textwidth}p{0.55\textwidth}@{}}
\toprule
\textbf{Idea conversada} \ideageo & \textbf{Dónde encaja} \\
\midrule
Recuperar la presión en profundidad &
Familia \textbf{A1}: estimación de $\Delta P$ y, con ella, de $P(z)$.
\alert{Hay que precisar el enunciado} (lámina siguiente). \\[0.4em]
Contenido de burbujas en profundidad &
Familia \textbf{B3}: recuperar el perfil $\phi(z,0)$. Problema de
observabilidad a lo largo de características. \\[0.4em]
Modelo más realista &
Familia \textbf{E1/E4}: unificar solvers, exsolución y cristalización en
desequilibrio, Darcy vertical y lateral. \\[0.4em]
Mostrar la no-unicidad respecto de la geometría &
Familia \textbf{D1}: \alert{ya demostrado y verificado} (sección 4). \\[0.4em]
Complejizar la geometría &
Familia \textbf{B1/E7}: $R(z)$ como \alert{incógnita funcional}, no sólo como
estudio paramétrico. \\
\bottomrule
\end{tabular}

\mensaje{Las cinco ideas son buenas y encajan. Lo que agrego es una
\textbf{columna vertebral} que las conecta —la no-unicidad y la
identificabilidad— y unas cuantas ideas nuevas.}
\end{frame}

% ---------------------------------------------------------------------------
\begin{frame}{Una precisión sobre ``recuperar la presión en profundidad''}
\begin{alertblock}{El enunciado ingenuo no está bien planteado}
$P(z)$ \textbf{no} es un dato independiente: una vez fijados los parámetros, el
modelo directo la produce. Lo que de verdad se invierte es la
\alert{condición de borde} $P(H)$ y, con ella, el perfil completo.
\end{alertblock}
\medskip
Y hay una mala noticia empírica que veremos en la sección 4: $\Delta P$ es
justamente el parámetro \alert{peor determinado} por los observables de
superficie.
\medskip
\begin{exampleblock}{El enunciado defendible}
``¿Qué \textbf{cota superior} sobre la sobrepresión de cámara imponen los
observables de la erupción, y \textbf{qué dato adicional} haría falta para
acotarla por abajo?''
\end{exampleblock}
\medskip
Esto no es una rebaja del objetivo: es el objetivo bien formulado. Y la
respuesta —``hace falta deformación o gas''— es un resultado útil para el
monitoreo.
\end{frame}

% ===========================================================================
\section{Marco matemático}
% ===========================================================================

\begin{frame}{El operador directo}
$$\mathcal{F}:\ \Theta\longrightarrow\mathcal{D},\qquad \theta\longmapsto d=\mathcal{F}(\theta)$$

\begin{columns}[T]
\begin{column}{0.48\textwidth}
\textbf{$\Theta$ de dimensión finita}
\begin{itemize}\small
  \item radio $R$, sobrepresión $\Delta P$
  \item agua inicial $c_0$, temperatura $T$
  \item cristales $\xi$, umbral $\phi_{crit}$
  \item constantes de permeabilidad
\end{itemize}
\end{column}
\begin{column}{0.48\textwidth}
\textbf{$\Theta$ de dimensión infinita}
\begin{itemize}\small
  \item la geometría $R(\cdot)$
  \item un perfil de volátiles $c_0(\cdot)$
  \item fricción de pared $\beta(\cdot)$
  \item historia de cámara $P_{cam}(\cdot)$
\end{itemize}
\end{column}
\end{columns}

\vspace{1em}
\begin{columns}[T]
\begin{column}{0.48\textwidth}
\textbf{Determinista}
$$\theta_\alpha\in\arg\min_\theta\ \tfrac12\|\mathcal{F}(\theta)-d^{obs}\|^2_{\Sigma^{-1}}+\alpha\,\mathcal{R}(\theta)$$
\footnotesize $\mathcal{R}$: Tikhonov, variación total (geometrías con saltos)
o Sobolev $\|R'\|^2_{L^2}$ (geometrías suaves). $\alpha$ por discrepancia de
Morozov o validación cruzada.
\end{column}
\begin{column}{0.48\textwidth}
\textbf{Bayesiano}
$$\pi(\theta\mid d^{obs})\propto e^{-\frac12\|\mathcal{F}(\theta)-d^{obs}\|^2_{\Sigma^{-1}}}\ \pi_0(\theta)$$
\footnotesize Adecuado aquí porque la petrología entrega \alert{previas reales}
(inclusiones vítreas $\to c_0$; equilibrios de fases $\to T$, profundidad) y
porque el $\MER$ se conoce \alert{salvo un factor 2}: una estimación puntual sin
barras de error no significa nada.
\end{column}
\end{columns}
\end{frame}

% ---------------------------------------------------------------------------
\begin{frame}{Por qué ``agregar dimensiones'' es lo que hace posible invertir}
\begin{itemize}
  \item<1-> \textbf{Estacionario 1D} $\rightarrow$ un puñado de \emph{escalares}
    por simulación. Pocos datos, muchos parámetros: masivamente subdeterminado.
  \item<2-> \textbf{Transiente} $\rightarrow$ una \emph{función del tiempo}.
    Cada instante es una restricción adicional, y el sistema recorre un tramo de
    la curva característica $\MER_{ss}(P_{cam})$, informando sobre su
    \alert{pendiente}, que es una función distinta de los parámetros.
  \item<3-> \textbf{2D} $\rightarrow$ un \emph{perfil radial}. La reología
    dependiente de la tasa de corte deja firma en la distribución radial de
    deformación, ligada a la sismicidad de borde de conducto y a las texturas de
    piroclastos de borde frente a centro.
\end{itemize}

\onslide<4->{
\mensaje{En lenguaje de problemas inversos: cada dimensión añadida
\textbf{aumenta el rango efectivo del operador linealizado} y por tanto reduce
la dimensión del espacio nulo. \textbf{Cuantificar esa reducción} —con la matriz
de información de Fisher y con perfiles de verosimilitud— es un resultado
concreto y demostrable, y es lo que conecta el trabajo previo con la tesis.}}
\end{frame}

% ===========================================================================
\section{Resultados preliminares}
% ===========================================================================

\begin{frame}{El prototipo: un recorte, no un reemplazo}
El modelo de la tesis es el \alert{bifásico} (\texttt{RIconduitex5\_5.py}:
$u_m\neq u_g$, arrastre, cristalización). El prototipo
(\texttt{tesis/prototipo/}) es un \emph{recorte} monofásico del mismo sistema,
con las mismas constitutivas, hecho para demostrar teoremas:
\medskip
\begin{columns}[T]
\begin{column}{0.5\textwidth}
\begin{itemize}\small
  \item[(i)] \textbf{barato} — estacionario en $\sim0.1$ s;
  \item[(ii)] \textbf{suave} — sin \emph{switches} ni bisecciones anidadas;
\end{itemize}
\end{column}
\begin{column}{0.5\textwidth}
\begin{itemize}\small
  \item[(iii)] \textbf{abierto a $R(z)$ arbitrario};
  \item[(iv)] consistente con el bifásico en el tramo viscoso
    pre-fragmentación.
\end{itemize}
\end{column}
\end{columns}
\medskip
Lubricación compresible en $z\in[H,z_f]$, con exsolución en equilibrio:
$$\frac{dP}{dz}=-\rho(P)\,g-\frac{8\,\mu(P)\,\MER}{\pi\,\rho(P)\,R(z)^4},
\qquad P(H)=P_{lit}+\Delta P,\quad P(z_f)=P_{frag}.$$

\mensaje{Caso base de Calbuco ($R=16$ m, $\Delta P=5$ MPa, 4 wt\% H$_2$O,
970 °C, 25 \% cristales): $\MER=3.0\times10^{7}$ kg/s, del mismo orden que los
$1.4$–$2.7\times10^{7}$ kg/s observados. \textbf{El caso base es
cuantitativamente razonable.}}
\end{frame}

% ---------------------------------------------------------------------------
\begin{frame}{Resultado 1 — Teorema de degeneración geométrica}
\small
En el régimen viscoso dominado por fricción, el balance de momentum con
conservación de masa se integra exactamente:
\vspace{-0.4em}
$$\int_{P_{frag}}^{P_{cam}}\frac{\rho(P)}{\mu(P)}\dd P
=\frac{8\,\MER}{\pi}\,J_4[R],\qquad
\boxed{\ J_4[R]=\int_H^{z_f}R(z)^{-4}\dd z\ }$$
\vspace{-1.0em}

\begin{alertblock}{El lado izquierdo no depende de la geometría}
Por lo tanto $\MER\cdot J_4[R]$ es \textbf{invariante}: dos conductos con el
mismo $J_4$ son \alert{exactamente} indistinguibles a partir de la tasa de
descarga estacionaria.
\end{alertblock}

La derivada de Fréchet es $DJ_4[R]\,\delta R=-4\int R^{-5}\delta R\dd z$, de
modo que el espacio nulo del operador linealizado es
\vspace{-0.4em}
$$\mathcal{N}=\Big\{\delta R\ :\ \int_H^{z_f}R(z)^{-5}\,\delta R(z)\dd z=0\Big\},$$
\vspace{-1.2em}

un subespacio \alert{cerrado de codimensión 1}: infinitas direcciones de
perturbación de la geometría son invisibles.

\mensaje{Verificación numérica: $\MER\cdot J_4$ constante a $<10^{-6}$ relativo
entre cinco formas distintas; la fórmula analítica reproduce el $\MER$ numérico
con error $+0.000\pc$.}
\end{frame}

% ---------------------------------------------------------------------------
\begin{frame}{Resultado 1 — la no-unicidad es enorme, no marginal}
\centering
\fig[0.98\linewidth]{exp1_no_unicidad_geometria.png}

\vspace{-0.2em}
\mensaje{Seis geometrías cualitativamente distintas, \textbf{con volúmenes que
difieren en $+317\pc$}, producen el mismo $\MER$ con error relativo
$\sim10^{-7}$ en el modelo \textbf{completo} (flotabilidad, exsolución,
viscosidad de Giordano, cristales).}
\end{frame}

% ---------------------------------------------------------------------------
\begin{frame}{Resultado 1 — la tabla}
\centering\small
\begin{tabular}{@{}lrrrrrr@{}}
\toprule
\textbf{Geometría} & $R(H)$ & $R(z_f)$ & \textbf{Volumen} & $\Delta V$ & $J_4$ & \textbf{error MER} \\
\midrule
cilindro                    & 16.0 & 16.0 & $4.83\times10^6$ & —        & $9.16\times10^{-2}$ & $+0.00000\pc$ \\
ensanche basal $\times3$     & 44.4 & 14.8 & $1.29\times10^7$ & $+166\pc$ & $5.86\times10^{-2}$ & $-0.00001\pc$ \\
ensanche basal $\times4$     & 51.4 & 13.6 & $2.01\times10^7$ & \alert{$+317\pc$} & $4.24\times10^{-2}$ & $-0.00002\pc$ \\
ensanchamiento superior $\times2$ & 11.2 & 22.4 & $4.74\times10^6$ & $-1.8\pc$ & $2.18\times10^{-1}$ & $+0.00000\pc$ \\
dos tramos $\times2$ (tipo NC2) & 10.7 & 21.5 & $5.38\times10^6$ & $+11.5\pc$ & $2.22\times10^{-1}$ & $+0.00001\pc$ \\
constricción profunda $-35\pc$ & 16.8 & 16.9 & $4.80\times10^6$ & $-0.6\pc$ & $1.16\times10^{-1}$ & $+0.00001\pc$ \\
\bottomrule
\end{tabular}

\medskip
Nótese que $J_4$ varía en un \alert{factor 5} entre estas geometrías y aun así
el $\MER$ es el mismo: con flotabilidad el invariante es otro funcional, pero
sigue siendo \textbf{un solo funcional escalar}.
\end{frame}

% ---------------------------------------------------------------------------
\begin{frame}{Resultado 1 — ¿dónde ``mira'' el dato?}
\begin{columns}[c]
\begin{column}{0.58\textwidth}
\fig[\linewidth]{exp1_kernel_resistencia.png}
\end{column}
\begin{column}{0.40\textwidth}
La densidad de resistencia $w(z)\propto\mu(z)/R(z)^4$ se concentra brutalmente
hacia arriba, porque la desgasificación endurece el magma varios órdenes de
magnitud:
\medskip
\begin{itemize}
  \item 50 \% de la resistencia en los últimos \alert{830 m};
  \item 90 \% en los últimos \alert{44 m};
  \item 99 \% en los últimos \alert{2 m}.
\end{itemize}
\end{column}
\end{columns}

\mensaje{\textbf{Consecuencia geológica directa:} cualquier inferencia sobre el
conducto \emph{profundo} hecha a partir de la tasa de descarga es, esencialmente,
una inferencia sobre \textbf{la previa} y no sobre \textbf{el dato}. Esto es una
advertencia metodológica publicable por sí sola.}
\end{frame}

% ---------------------------------------------------------------------------
\begin{frame}{Resultado 2 — el estacionario confunde $R$ con $\Delta P$}
\begin{columns}[c]
\begin{column}{0.46\textwidth}
\fig[\linewidth]{exp2_degeneracion_RdP.png}
\end{column}
\begin{column}{0.52\textwidth}
El conjunto de nivel
$\{(R,\Delta P):\MER_{ss}=\MER_{obs}\}$ es una \alert{curva}: el dato
estacionario fija \textbf{una} combinación de los dos parámetros, no los dos.
\medskip

Matriz de sensibilidad en el punto nominal ($R=16$ m, $\Delta P=5$ MPa):
$$J=\frac{\partial(\log\MER,\log G)}{\partial(\log R,\log\Delta P)}
=\begin{pmatrix}4.000 & 0.195\\ 4.000 & 0.115\end{pmatrix}$$

\footnotesize
El $4.000$ es la ley de Poiseuille $\MER\propto R^4$ \emph{recuperada
numéricamente}. El $0.195$ es pequeño porque la sobrepresión aporta sólo una
fracción del empuje total (el resto lo aporta la flotabilidad): de ahí la ley de
compensación local $\Delta P\sim R^{-20.5}$.
\end{column}
\end{columns}

\mensaje{\textbf{Variaciones enormes de sobrepresión se compensan con
variaciones minúsculas de radio.}}
\end{frame}

% ---------------------------------------------------------------------------
\begin{frame}{Resultado 2 — el tiempo los separa}
\centering
\fig[0.96\linewidth]{exp2_curvas_MER_tiempo.png}

\vspace{-0.3em}
{\small Conductos con $\MER$ inicial \textbf{idéntico} ($3.01\times10^7$ kg/s)
drenando una misma cámara elástica de 50 km$^3$.}

\mensaje{Comparten un \emph{punto} de la curva característica
$\MER_{ss}(P_{cam})$, pero no su \alert{pendiente}
$G=d\MER/dP_{cam}\approx\pi R^4/(8\mu L)$, que fija el tiempo de decaimiento
$\tau=C/G$ de la erupción.}
\end{frame}

% ---------------------------------------------------------------------------
\begin{frame}{Resultado 2 — la tabla}
\centering\small
\begin{tabular}{@{}rrrrrr@{}}
\toprule
$R$ [m] & $\Delta P$ [MPa] & $G$ [kg/s/Pa] & $\tau$ [h] & $\MER$ a 6 h & masa en 6 h [kg] \\
\midrule
10 & 92.5 & 0.353 & 19.5 & $2.23\times10^7$ & $5.62\times10^{11}$ \\
12 & 45.3 & 0.617 & 11.1 & $1.84\times10^7$ & $5.10\times10^{11}$ \\
14 & 20.6 & 0.905 &  7.6 & $1.55\times10^7$ & $4.66\times10^{11}$ \\
16 &  5.0 & 1.174 &  5.8 & $1.36\times10^7$ & $4.34\times10^{11}$ \\
\bottomrule
\end{tabular}

\vspace{1.2em}
Los cuatro son \alert{indistinguibles} con un dato estacionario.
Con la serie de tiempo, los tiempos característicos van de \alert{5.8 h a
19.5 h} y las masas emitidas en seis horas difieren en \alert{29 \%}.

\mensaje{Ésta es la razón técnica por la que el trabajo previo —agregar el
tiempo al modelo— \textbf{no era un adorno}: es lo que vuelve identificable un
parámetro que antes no lo era.}
\end{frame}

% ---------------------------------------------------------------------------
\begin{frame}{Resultado 3 — cuánto se gana, cuantificado}
\centering
\figh[0.80\textheight]{exp3_identificabilidad.png}

\vspace{0.1em}
{\scriptsize Operador restringido
$\theta=(\log R,\log\Delta P)\mapsto d=(\log\MER,\log G)$; perfiles de
verosimilitud $\chi^2_{perfil}(\theta_i)=\min_{\theta_j,j\neq i}\chi^2(\theta)$
(Raue \emph{et al.}, 2009).}
\end{frame}

% ---------------------------------------------------------------------------
\begin{frame}{Resultado 3 — el veredicto}
Con un modelo de ruido realista ($\sigma_{\log\MER}=\ln2/2$, es decir un factor
2 de incertidumbre en la tasa de descarga):

\vspace{0.5em}
\centering\small
\begin{tabular}{@{}lcc@{}}
\toprule
\textbf{Datos} & \textbf{intervalo 95 \% de $R$} & \textbf{intervalo 95 \% de $\Delta P$} \\
\midrule
sólo estacionario & $[8,\,19]$ m (abierto) & $[0.2,\,250]$ MPa (todo el rango) \\
estacionario + serie de tiempo & $[13,\,18]$ m & $[0.25,\,44]$ MPa (sólo cota superior) \\
\bottomrule
\end{tabular}

\vspace{1.0em}
\begin{columns}[T]
\begin{column}{0.92\textwidth}
\begin{alertblock}{Conclusión honesta}
La serie de tiempo \textbf{identifica el radio} y \textbf{acota por arriba} la
sobrepresión, pero \textbf{no la determina}. Para determinarla hacen falta datos
de \alert{otra naturaleza}: deformación, emisión de gas.
\end{alertblock}
\end{column}
\end{columns}

\mensaje{Esto es \textbf{un resultado, no un fracaso}: es la justificación
cuantitativa de la inversión conjunta, y una respuesta concreta a la pregunta
``¿qué habría que medir?''.}
\end{frame}

% ---------------------------------------------------------------------------
\begin{frame}{Resultado 4 — ¿cuándo hace falta de verdad el $\partial_t$?}
\small
``Agregar el tiempo'' significa dos cosas distintas, y conviene no confundirlas:
\smallskip
\begin{columns}[T]
\begin{column}{0.47\textwidth}
\begin{block}{(A) El observable depende de $t$}
$\MER(t)$ es una función, no un escalar. Casi siempre ayuda a invertir
(resultado 2).
\end{block}
\end{column}
\begin{column}{0.47\textwidth}
\begin{block}{(B) El conducto está fuera de equilibrio}
Los términos $\partial_t$ de sus ecuaciones importan. Hay que resolver la EDP.
\end{block}
\end{column}
\end{columns}
\smallskip
\alert{(A) no implica (B).} El criterio es un número de Deborah
$\mathrm{De}=\tau_{conducto}/\tau_{forzamiento}$, con escalas \emph{medidas} del
propio modelo:

\vspace{0.3em}
\centering\footnotesize
\begin{tabular}{@{}lrl@{}}
\toprule
\textbf{Escala} & \textbf{Valor} & \textbf{Qué es} \\
\midrule
$\tau_{hidr}$  & 8 s     & relajación de $\MER$ tras un escalón de $-2\pc$ en $P_{cam}$ \\
$\tau_{resid}$ & 6.0 min & masa en el conducto dividida por $\MER$ \\
$\tau_{cam}$   & 5.83 h  & drenaje de la cámara, $C/G$, con $V=50$ km$^3$ \\
\bottomrule
\end{tabular}

\mensaje{$\mathrm{De}=5.7\times10^{-4}$: el conducto está \textbf{esclavizado} a
la cámara. Alimentando el transiente completo con el $P_{cam}(t)$
cuasi-estacionario, el error en $\MER(t)$ es de \textbf{0.08 \% máximo}.}
\end{frame}

% ---------------------------------------------------------------------------
\begin{frame}{Resultado 4 — la cinética, no la hidráulica, es la que decide}
\centering
\figh[0.52\textheight]{exp4_escalas_de_tiempo.png}

\vspace{0.2em}
\begin{columns}[T]
\begin{column}{0.95\textwidth}
\footnotesize
El $\tau$ medido es sólo el \textbf{hidráulico}: el modelo reducido supone
equilibrio, así que no contiene las escalas lentas que la literatura identifica
como las que de verdad vuelven transiente una erupción —exsolución en
desequilibrio ($10^1$–$10^3$ s), cristalización de microlitos ($10^4$–$10^6$ s),
escape de gas ($10^3$–$10^5$ s). Con $\tau\sim10^4$ s, hasta una fase
sub-pliniana de 6 h pasa a $\mathrm{De}\approx0.5$.
\end{column}
\end{columns}

\mensaje{\textbf{Consecuencia para la tesis: lo que se puede invertir depende
del régimen.} En cuasi-estacionario los datos informan sobre geometría y cámara
y \alert{nada} sobre cinética. En los ciclos de domo el período lo fija la
propia cinética, de modo que $\mathrm{De}\approx1$ por construcción: ahí el
transiente no es un refinamiento, \textbf{es el mecanismo}.}
\end{frame}

% ---------------------------------------------------------------------------
\begin{frame}{Síntesis de los resultados preliminares}
\begin{columns}[T]
\begin{column}{0.24\textwidth}
\begin{block}{1. Degeneración}
\footnotesize El $\MER$ estacionario ve $R(\cdot)$ sólo a través de \alert{un
escalar}. La preimagen de un dato es una hipersuperficie de codimensión 1 en un
espacio de funciones.\smallskip

\emph{Demostrado y verificado.}
\end{block}
\end{column}
\begin{column}{0.24\textwidth}
\begin{block}{2. Es enorme}
\footnotesize $+317\pc$ de diferencia de volumen con el mismo $\MER$. El 90 \%
de la caída de presión está en el kilómetro más somero: \alert{el conducto
profundo es invisible}.
\end{block}
\end{column}
\begin{column}{0.24\textwidth}
\begin{block}{3. El tiempo ayuda}
\footnotesize Pero no lo resuelve todo: identifica $R$, sólo acota $\Delta P$
por arriba. \alert{Cuánto aporta cada dimensión se puede medir.}
\end{block}
\end{column}
\begin{column}{0.24\textwidth}
\begin{block}{4. Y hay dos tiempos}
\footnotesize El observable dependiente de $t$ y el conducto fuera de
equilibrio son cosas distintas. \alert{El régimen decide qué es invertible.}
\end{block}
\end{column}
\end{columns}

\vspace{1.2em}
\mensaje{Los cuatro son reproducibles hoy: \texttt{tesis/prototipo/exp1..exp4}.
\textbf{La propuesta no parte de cero: parte de resultados.}}
\end{frame}

% ===========================================================================
\section{Catálogo de problemas inversos}
% ===========================================================================

\begin{frame}{Familia A — parámetros (dimensión finita)}
\small
\begin{tabular}{@{}p{0.05\textwidth}p{0.32\textwidth}p{0.26\textwidth}p{0.22\textwidth}@{}}
\toprule
\# & \textbf{Qué se recupera} & \textbf{Datos} & \textbf{Herramienta} \\
\midrule
\textbf{A1}\,\ideageo & Sobrepresión $\Delta P$ y, con ella, $P(z)$ &
$\MER$, duración, masa, altura de columna & MCMC con previas petrológicas \\[0.4em]
\textbf{A2}\,\ideageo & Volátiles $c_0$ y cristales $\xi$ &
$+$ inclusiones vítreas & Previa informativa fuerte \\[0.4em]
\textbf{A3} & Profundidad de fragmentación $z_f$, umbral $\phi_{crit}$ &
Densidad, porosidad y vesicularidad de piroclastos & Inversión conjunta con A1–A2 \\[0.4em]
\textbf{A4} & Constantes de $k(\phi)$ &
Porosidad–permeabilidad de laboratorio $+$ $\MER$ & Datos de dos fuentes distintas \\
\bottomrule
\end{tabular}

\mensaje{\ideageo{} $=$ ideas ya conversadas con el profesor de geología.}
\end{frame}

% ---------------------------------------------------------------------------
\begin{frame}{Familia B — funciones (dimensión infinita)}
\footnotesize
\begin{tabular}{@{}p{0.05\textwidth}p{0.33\textwidth}p{0.25\textwidth}p{0.23\textwidth}@{}}
\toprule
\# & \textbf{Qué se recupera} & \textbf{Datos} & \textbf{Herramienta} \\
\midrule
\textbf{B1}\,\ideageo\ideageo & \alert{Geometría $R(z)$} &
$\MER(t)$, presión de salida, $z_f$ & Tikhonov / TV, con el teorema de $J_4$
como resultado de no-unicidad \\[0.4em]
\textbf{B2} & Fricción de pared $\beta(z)$, condición de Navier
$\mu\,\partial_r u=-\beta u$ en $r=R$ &
$\MER(t)$, sismicidad de borde, texturas de borde &
Coeficiente de Robin: unicidad y estabilidad logarítmica conocidas \\[0.4em]
\textbf{B3}\,\ideageo & Porosidad inicial $\phi(z,0)$ &
Serie de porosidad/densidad de piroclastos en el vent &
Observabilidad sobre características \\[0.4em]
\textbf{B4} & Ley $k(\phi)$ como \emph{función} & igual que A4 &
Identificación de coeficiente no lineal \\
\bottomrule
\end{tabular}

\mensaje{\textbf{B1 es el núcleo de la tesis.} Y \textbf{B2 es la extensión
natural del código 2D que ya existe}: \texttt{RIconduit2D\_FD.py} impone
no-deslizamiento $u_m(R)=0$; reemplazarlo por Navier con $\beta(z)$ desconocido
convierte el problema en identificación de un coeficiente de Robin desde la
frontera —con teoría establecida. Geológicamente $\beta(z)$ es la zona de
cizalle del borde, asociada a los sismos repetitivos de Mount St. Helens.}
\end{frame}

% ---------------------------------------------------------------------------
\begin{frame}{Familia C — historias temporales}
\footnotesize
\begin{tabular}{@{}p{0.05\textwidth}p{0.30\textwidth}p{0.25\textwidth}p{0.29\textwidth}@{}}
\toprule
\# & \textbf{Qué se recupera} & \textbf{Datos} & \textbf{Herramienta} \\
\midrule
\textbf{C1} & Historia de presión de cámara $P_{cam}(t)$ &
$\MER(t)$ de la columna eruptiva &
Problema lateral (\emph{sideways}) / deconvolución de Volterra, mal puesto \\[0.4em]
\textbf{C2} & Historia de recarga $Q_{in}(t)$ &
$\MER(t)$ $+$ deformación & Igual, con cámara–conducto acoplado \\[0.4em]
\textbf{C3} & Tasa de descompresión $dP/dt$ durante el ascenso &
Densidad numérica de microlitos en piroclastos &
Inversión de una ecuación de balance poblacional \\
\bottomrule
\end{tabular}

\vspace{0.6em}
\scriptsize
\textbf{C1 es atractivo por su limpieza matemática.} La ecuación parabólica no
lineal del prototipo
$A(z)\rho'(P)\partial_tP=\partial_z[k(z)\rho(P)(\partial_zP+\rho g)]$, con dato
$\MER(t)$ en $z=z_f$ y la incógnita en $z=H$, es un problema de frontera lateral
para una ecuación parabólica: severamente mal puesto, con teoría clásica de
estabilidad condicional. \emph{Es un tema de tesis de ingeniería matemática por
derecho propio.}

\smallskip
\textbf{C3 conecta directamente con la geología de terreno.} La densidad
numérica de microlitos es un ``velocímetro'' de descompresión ampliamente usado.
Formularlo como problema inverso explícito —con regularización y barras de
error, en vez de una calibración empírica— sería una contribución metodológica
concreta a la petrología.
\end{frame}

% ---------------------------------------------------------------------------
\begin{frame}{Familia D — teoría de unicidad y no-unicidad}
\small
\begin{tabular}{@{}p{0.05\textwidth}p{0.45\textwidth}p{0.30\textwidth}p{0.10\textwidth}@{}}
\toprule
\# & \textbf{Resultado buscado} & \textbf{Herramienta} & \textbf{Estado} \\
\midrule
\textbf{D1}\,\ideageo & Teorema de degeneración $J_4$ y caracterización del
espacio nulo & Cálculo variacional & \alert{hecho} \\[0.4em]
\textbf{D2}\,\ideageo & No-unicidad de la geometría desde datos sísmicos de
resonancia & Ecuación de Webster $+$ teoría espectral inversa de
Sturm–Liouville & propuesto \\[0.4em]
\textbf{D3} & Multiplicidad de estados estacionarios (curva sigmoidal
$p_{ch}$–$q$) como fuente \emph{adicional} de no-unicidad &
Teoría de bifurcaciones; Melnik \& Sparks (1999), Kozono \& Koyaguchi (2012) &
propuesto \\
\bottomrule
\end{tabular}

\mensaje{Ojo con la distinción: \textbf{D3} es la no-unicidad \emph{del problema
directo} (varios caudales para los mismos parámetros), que la literatura sí
conoce. \textbf{D1 y D2} son no-unicidad \emph{del problema inverso}
(no-inyectividad de $\mathcal{F}$), que —hasta donde he podido revisar— no ha
sido formalizada en volcanología.}
\end{frame}

% ---------------------------------------------------------------------------
\begin{frame}{D2: la joya matemática de la propuesta}
\small
Linealizando el flujo compresible en un conducto de sección $A(z)$ se obtiene la
\textbf{ecuación de Webster}, que con $p=A^{-1/2}\psi$ se transforma en un
problema de \textbf{Sturm–Liouville}:
\vspace{-0.3em}
$$\frac{\partial^2p}{\partial t^2}=c^2\,\frac1A\frac{\partial}{\partial z}\!\left(A\frac{\partial p}{\partial z}\right)
\qquad\Longleftrightarrow\qquad
\psi''+\left(\frac{\omega^2}{c^2}-V(z)\right)\psi=0,\quad V=\frac{(A^{1/2})''}{A^{1/2}}.$$
\vspace{-1.0em}

Dos hechos clásicos con lectura volcanológica inmediata:
\begin{itemize}
  \item \textbf{Una sola familia espectral no determina el potencial} (teorema
    de dos espectros de Borg, 1946). Es decir: \alert{el espectro de resonancia
    de un conducto no determina su geometría}.
  \item Aun conociendo $V$ exactamente, $(A^{1/2})''=V\,A^{1/2}$ tiene un
    espacio de soluciones de \alert{dimensión 2}: existe una \textbf{familia
    biparamétrica explícita} de secciones $A(z)$ acústicamente indistinguibles.
\end{itemize}

\mensaje{Como en volcanología se infieren rutinariamente dimensiones de conducto
a partir de eventos de largo período y tremor, exhibir esa familia con
parámetros realistas sería una \textbf{advertencia metodológica fuerte y
matemáticamente irreprochable}.}
\end{frame}

% ---------------------------------------------------------------------------
\begin{frame}{Familia E — metodología y realismo del modelo}
\footnotesize
\begin{tabular}{@{}p{0.05\textwidth}p{0.43\textwidth}p{0.43\textwidth}@{}}
\toprule
\# & \textbf{Qué se hace} & \textbf{Por qué importa} \\
\midrule
\textbf{E1}\,\ideageo & Unificar los tres solvers en \alert{un} operador
directo, con fragmentación regularizada &
Sin esto no hay gradientes ni inversión eficiente. \alert{Prerrequisito} \\[0.3em]
\textbf{E2} & Método \textbf{adjunto} para el sistema transiente &
Gradiente a costo $O(1)$ \emph{solves} en vez de $O(p)$: habilita $R(z)$ con
cientos de incógnitas \\[0.3em]
\textbf{E3} & \textbf{Selección de modelo} entre criterios de fragmentación &
Convierte ``¿cuál criterio es correcto?'' en una pregunta con respuesta
cuantitativa (factores de Bayes) \\[0.3em]
\textbf{E4}\,\ideageo & Modelo más realista: desequilibrio de exsolución y
cristalización, Darcy vertical y lateral, cámara acoplada &
Wong \emph{et al.} (2017): la permeabilidad controla el flujo másico
erupcionado \\[0.3em]
\textbf{E5} & \textbf{Emulador} (proceso gaussiano o red neuronal) &
Hace factible MCMC con $10^5$–$10^6$ evaluaciones \\[0.3em]
\textbf{E6} & \textbf{Diseño óptimo de experimentos} &
Qué observable reduce más la incertidumbre: útil para el monitoreo
volcánico en Chile \\[0.3em]
\textbf{E7}\,\ideageo & Geometría más compleja: dique profundo $\to$ cilindro
somero, sección elíptica, cráter &
Aravena \emph{et al.} (2017, 2018): efectos de primer orden en presión de salida
y $\MER$ \\
\bottomrule
\end{tabular}
\end{frame}

% ===========================================================================
\section{Estado del arte}
% ===========================================================================

\begin{frame}{Sí existen, y son bastantes: modelos con TIEMPO}
\footnotesize
\begin{itemize}
  \item \textbf{Melnik \& Sparks (1999)}, \emph{Nature} 402:37 — el resultado
    fundacional: para condiciones fijas puede haber \alert{hasta tres tasas de
    descarga estacionarias}, que difieren en órdenes de magnitud.
  \item \textbf{Barmin, Melnik \& Sparks (2002)}, \emph{EPSL} 199:173 — cámara
    $+$ conducto reducido que produce ciclos por transiciones entre ramas
    estables. Matemáticamente limpio: buen modelo de juguete para probar
    algoritmos de inversión.
  \item \textbf{Melnik \& Sparks (2005)}, \emph{JGR} — periodos controlados
    principalmente por el \alert{volumen de la cámara}. Confirma el mecanismo
    que usa nuestro experimento 2.
  \item \textbf{La Spina, de' Michieli Vitturi \& Clarke (2017)}, \emph{JVGR}
    336:118 — \alert{lo más cercano a nuestro \texttt{RIconduit1D\_transient},
    pero más completo}: sistema hiperbólico con dos presiones y dos velocidades,
    tasas finitas de exsolución y relajación, flujo denso y diluido unificados.
    Conclusión: \emph{la duración de la erupción la controlan las escalas de
    tiempo de exsolución}, no el equilibrio. \textbf{Lectura obligada antes de
    fijar el operador directo.}
  \item \textbf{Wong \& Segall (2019)}, \emph{JGR} — transiente con escape de gas
    de Darcy vertical y lateral. Observación relevante: \alert{la mayoría de los
    transientes a escala de erupción siguen siendo 1D promediados
    radialmente}.
  \item \textbf{Código abierto \textsc{mamma}} (de' Michieli Vitturi, La Spina,
    Aravena) — vale la pena evaluarlo como operador directo alternativo, en vez
    de reescribir todo.
\end{itemize}
\end{frame}

% ---------------------------------------------------------------------------
\begin{frame}{Modelos con GEOMETRÍA VARIABLE}
\footnotesize
\begin{itemize}
  \item \textbf{de' Michieli Vitturi \emph{et al.} (2008)}, \emph{EPSL}
    272:567 — código \textsc{domeflow}, transiente ``1.5D'' con radio variable.
    Barren $0.4\le R_t/R_b\le2.5$ y $0.1\le H/L\le0.7$.
  \item \textbf{de' Michieli Vitturi \emph{et al.} (2010)}, \emph{EPSL} — más de
    300 corridas y análisis dimensional dan una relación universal: el producto
    del cambio de tasa de extrusión por el tiempo de transición es proporcional
    al cambio normalizado de presión de cámara, \alert{al volumen del conducto} y
    a la razón de radios. \emph{Es exactamente la degeneración que explota
    nuestro experimento 2}, y da un punto de comparación directo.
  \item \textbf{Costa, Melnik \& Sparks (2007)}, \emph{EPSL} 260:137 — acopla la
    geometría a la elasticidad del encajante: el conducto se deforma con la
    presión. Antecedente natural de un problema de frontera libre.
  \item \textbf{Aravena \emph{et al.} (2017, 2018)}, \emph{JVGR} / \emph{Sci.
    Rep.} — criterios de colapso de Mohr–Coulomb aplicados al perfil $P(z)$:
    \alert{los conductos cilíndricos son estables sólo para radios grandes}, de
    modo que el ensanchamiento sineruptivo lleva naturalmente a $R(z)$. Definen
    las configuraciones NC2 y NC3, que nuestro experimento 1 reproduce.
\end{itemize}

\mensaje{\textbf{Álvaro Aravena es chileno} y trabaja en esta línea exacta
(geometría no cilíndrica $+$ modelos de conducto). Contacto natural y posible
evaluador externo.}
\end{frame}

% ---------------------------------------------------------------------------
\begin{frame}{Modelos 2D}
\footnotesize
\begin{itemize}
  \item \textbf{Massol, Jaupart \& Pepper (2001)}, \emph{JGR} 106:16223 — el 2D
    axisimétrico de referencia. Resultados que un 1D no puede dar: la
    sobrepresión de gas en el vent supera 1 MPa; hay variaciones
    \alert{horizontales} grandes de presión de gas y agua exsuelta; y el
    resultado depende críticamente de la \alert{condición de borde de salida}
    (velocidad horizontal nula vs.\ esfuerzo de corte nulo cambian la
    sobrepresión en $\sim$1 MPa).
  \item \textbf{Collier \& Neuberg (2006)}, \textbf{Thomas \& Neuberg (2014,
    2019)} — 2D axisimétrico por elementos finitos con tres fases, pérdida de
    gas, deslizamiento en la pared y falla frágil; el de 2019 acopla el flujo con
    deformación elástica para predecir inclinometría. \emph{Precedente directo de
    una inversión conjunta flujo $+$ geodesia.}
  \item \textbf{Dufek \& Bergantz (2005)}, \textbf{Tsvetkova \& Melnik (2018)},
    \textbf{Longo \emph{et al.} (\textsc{gales})} — transiente y 2D
    simultáneamente, reología dependiente de la tasa de corte, multicomponente.
\end{itemize}

\mensaje{El punto de Massol \emph{et al.} sobre la condición de borde de salida
es directamente relevante para la propuesta: es un ejemplo documentado de que
\textbf{la condición de borde importa tanto como los parámetros}, lo que
justifica tratarla como incógnita (idea B2).}
\end{frame}

% ---------------------------------------------------------------------------
\begin{frame}{Modelos 3D: la recomendación es no}
\begin{alertblock}{Observación honesta}
Los modelos genuinamente 3D de conducto son escasos, \emph{y por buena razón
física}: el conducto es un tubo esbelto ($L/R\sim400$ en Calbuco), de modo que la
axisimetría captura casi toda la física. El 3D se justifica sólo para diques,
secciones elípticas, conductos inclinados o sistemas múltiples. El retorno
científico es bajo frente al costo, y el problema inverso se vuelve intratable.
\alert{La dimensión que sí paga es el tiempo, y después la geometría variable.}
\end{alertblock}

\medskip
Dicho de otro modo: el costo computacional de una dimensión extra sólo se
justifica si esa dimensión \emph{aporta información al problema inverso}. El
experimento 3 muestra que el tiempo sí la aporta, y de manera medible. Para el
3D, en un tubo con $L/R\sim400$, no hay razón para esperar lo mismo.

\medskip
\small Donde el 3D sí sería necesario: un \textbf{dique} en profundidad que pasa
a conducto cilíndrico somero, una \textbf{sección elíptica} —ambos, de hecho,
casos de ``complejizar la geometría'' que se pueden tratar de forma efectiva en
un modelo 1.5D con un radio hidráulico equivalente (idea E7).
\end{frame}

% ---------------------------------------------------------------------------
\begin{frame}{Problemas inversos en volcanología}
La línea está esencialmente monopolizada por el grupo de Stanford/USGS.
\smallskip
\small
\begin{itemize}
  \item \textbf{Anderson \& Segall (2011, 2013)}, \emph{JGR} — la referencia
    fundacional. Invierten GPS $+$ tasa de extrusión con un modelo físico de
    cámara elipsoidal $+$ conducto cilíndrico, por MCMC. Argumento central: a
    diferencia de las inversiones geodésicas con modelos \emph{cinemáticos}, esta
    técnica \alert{sí} puede restringir propiedades del magma y el volumen y
    presión absolutos de la cámara.
  \item \textbf{Wong, Segall, Bradley \& Anderson (2017)}, \emph{JGR} 122:7789 —
    agregan cristalización en equilibrio y Darcy. Hallazgos relevantes: el flujo
    másico \alert{depende fuertemente de las permeabilidades}; incluir transporte
    de gas lateral \emph{y} vertical produce comportamiento \alert{no monótono}
    —es decir, no-unicidad adicional.
  \item \textbf{Wong \& Segall \emph{et al.} (2020)}, \emph{G-Cubed} — inversión
    conjunta de \alert{tres series de tiempo} (volumen extruido, deformación,
    CO$_2$). Es el estándar metodológico al que debería apuntar el capítulo 4.
\end{itemize}
\end{frame}

% ---------------------------------------------------------------------------
\begin{frame}{Lo que ese cuerpo de trabajo NO hace}
\begin{enumerate}
  \item<1-> \textbf{Ninguno trata la geometría del conducto como función
    incógnita.} Todos invierten un \emph{radio escalar} (a lo sumo un radio y una
    longitud de tapón).
  \item<2-> \textbf{Ninguno presenta un análisis de no-unicidad estructural.}
    Reportan posteriores anchas, pero no demuestran \emph{por qué} son anchas ni
    caracterizan el espacio nulo.
  \item<3-> \textbf{Todos son casos efusivos} (domos de Mount St. Helens,
    Soufrière Hills). No hay inversión bayesiana con modelo físico de una
    erupción \alert{explosiva sub-pliniana}.
  \item<4-> \textbf{Ninguno usa un modelo con dimensión radial resuelta.}
\end{enumerate}

\onslide<5->{
\mensaje{Ahí está el nicho —y las cuatro cosas son justamente lo que este
proyecto ya tiene a medio camino.}}
\end{frame}

% ---------------------------------------------------------------------------
\begin{frame}{Dónde encaja esta tesis}
\centering\footnotesize
\begin{tabular}{@{}p{0.18\textwidth}p{0.36\textwidth}p{0.39\textwidth}@{}}
\toprule
\textbf{Dimensión} & \textbf{Estado del arte} & \textbf{Qué falta} \\
\midrule
\textbf{Tiempo} & Maduro (Melnik \& Sparks; La Spina; Wong \& Segall) &
Usarlo \alert{para invertir}, no sólo para simular, y cuantificar cuánta
información aporta \\[0.4em]
\textbf{Radial (2D)} & Maduro (Massol \& Jaupart; Collier \& Neuberg) &
Nadie lo ha usado en un problema inverso; \alert{la condición de borde de pared
como incógnita es territorio virgen} \\[0.4em]
\textbf{3D} & Escaso y poco justificado físicamente & No recomendado \\[0.4em]
\textbf{Geometría variable} & Maduro como estudio paramétrico \emph{directo}
(de' Michieli Vitturi; Aravena; Costa) &
Nadie la ha tratado como \alert{incógnita de un problema inverso}, ni ha
caracterizado su no-unicidad \\[0.4em]
\textbf{Inversión bayesiana} & Establecida para casos efusivos &
Ningún caso explosivo sub-pliniano; ningún análisis estructural de no-unicidad;
ningún caso chileno \\
\bottomrule
\end{tabular}

\mensaje{\textbf{El nicho es la intersección de las dos últimas filas:} tratar la
geometría del conducto como incógnita funcional en un problema inverso
bayesiano, demostrar que es estructuralmente no único, caracterizar el espacio
nulo y determinar qué datos lo reducen — aplicado a una erupción explosiva
chilena con datos reales.}
\end{frame}

% ===========================================================================
\section{La tesis propuesta}
% ===========================================================================

\begin{frame}{Título y estructura}
\begin{block}{}
\centering
\textbf{``¿Qué se puede recuperar de una erupción? Identificabilidad y
no-unicidad en modelos de ascenso de magma por un conducto, con aplicación a
Calbuco 2015.''}
\end{block}

\vspace{0.4em}
\footnotesize
\begin{description}
  \item[\textbf{Cap.\ 1 — El operador directo} (E1)] Unificación de los tres
    solvers; regularización de la fragmentación; verificación de que la versión
    suave reproduce la original; costo y diferenciabilidad.
  \item[\textbf{Cap.\ 2 — Teoría de no-unicidad} (D1, D2)] El teorema $J_4$ y el
    espacio nulo. Webster–Sturm–Liouville y la familia biparamétrica de
    conductos acústicamente idénticos.
  \item[\textbf{Cap.\ 3 — Identificabilidad del bifásico} ] ¿La degeneración
    $J_4$ sobrevive a $u_g-u_m$? Fisher y perfiles sobre
    \texttt{RIconduitex5\_5.py}. Estacionario vs.\ transiente vs.\ 2D.
    \alert{Conecta el trabajo previo con la tesis.}
  \item[\textbf{Cap.\ 4 — Inversión bayesiana de Calbuco} (A1–A3)] Previas de
    Castruccio \emph{et al.} (2016); posterior sobre $\Delta P$, $R$, $c_0$,
    $z_f$; verificación predictiva posterior.
  \item[\textbf{Cap.\ 5 — Geometría como incógnita funcional} (B1, E7)]
    Inversión de $R(z)$ con regularización; demostración numérica de que el
    mínimo no es único; conexión con el ensanchamiento sineruptivo.
\end{description}

\mensaje{Si el tiempo alcanza, \textbf{B2} extiende el capítulo 5 hacia el
modelo 2D y da material para un artículo independiente.}
\end{frame}

% ---------------------------------------------------------------------------
\begin{frame}{Datos disponibles para Calbuco 2015}
\centering\scriptsize
\begin{tabular}{@{}p{0.31\textwidth}p{0.33\textwidth}p{0.30\textwidth}@{}}
\toprule
\textbf{Observable} & \textbf{Valor} & \textbf{Fuente} \\
\midrule
Volumen emitido total & 0.27–0.56 km$^3$ (no DRE) & Romero (2016); Castruccio (2016); Van Eaton (2016) \\
Reparto entre fases & $\sim$38 \% fase 1, $\sim$62 \% fase 2 & Castruccio (2016) \\
Masa fase 1 / fase 2 & $8\times10^{10}$ kg en 1.5 h / $3.2\times10^{11}$ kg en 6 h & Castruccio (2016) \\
$\MER$ media & $1.4$ y $1.5\times10^7$ kg/s & Castruccio (2016) \\
$\MER$ máxima (isopletas) & $2.4$ y $2.7\times10^7$ kg/s & Castruccio (2016) \\
Altura de columna & 15–17 km sobre el cráter & SERNAGEOMIN; Van Eaton (2016) \\
Composición & andesita basáltica, 54–55 \% SiO$_2$ & Romero (2016) \\
Almacenamiento & 8–12 km, 900–950 °C, 5.5–6.5 wt\% H$_2$O & Morgado (2019); Arzilli (2019) \\
Texturas de piroclastos & CSD, BSD, densidad, porosidad, permeabilidad & Castruccio (AGU 2018) \\
Deformación & co-eruptiva detectada; pre-eruptiva bajo el ruido & Delgado (2017) \\
Pluma y descargas eléctricas & \alert{serie temporal} & Van Eaton (2016) \\
\bottomrule
\end{tabular}

\mensaje{\textbf{Advertencia de consistencia:} \texttt{calbuco2015d.py} usa
$T=970$ °C, $c_0=4$ wt\% y base a 7 km; la literatura reciente sugiere
900–950 °C, 5.5–6.5 wt\% y 8–12 km. La viscosidad de Giordano es
\alert{exponencialmente} sensible a ambos. Conviene reconciliarlo antes de fijar
previas.}
\end{frame}

% ---------------------------------------------------------------------------
\begin{frame}{Una oportunidad concreta para el profesor de geología}
\begin{columns}[T]
\begin{column}{0.48\textwidth}
\begin{block}{Arzilli \emph{et al.} (2019)}
Gatillo \alert{interno}: cristalización prolongada $\to$ segunda ebullición
$\to$ sobrepresurización. \textbf{Sin} intrusión de magma fresco.
\end{block}
\end{column}
\begin{column}{0.48\textwidth}
\begin{block}{Morgado \emph{et al.} (2019)}
Gatillo \alert{externo}: calentamiento localizado del reservorio por inyección
de magma caliente.
\end{block}
\end{column}
\end{columns}

\vspace{1em}
Son \textbf{dos hipótesis publicadas y contrapuestas} sobre el mecanismo
gatillante de Calbuco 2015. Ambas predicen \alert{historias de sobrepresión
distintas}, y por tanto \alert{curvas $\MER(t)$ distintas}.

\vspace{1em}
\begin{exampleblock}{Propuesta}
Discriminar entre ellas mediante \textbf{selección bayesiana de modelos}
(factores de Bayes). Sería un resultado con impacto en la comunidad
volcanológica chilena, y es exactamente el tipo de pregunta que un problema
inverso bien planteado \emph{sí} puede responder.
\end{exampleblock}
\end{frame}

% ---------------------------------------------------------------------------
\begin{frame}{Plan de trabajo y riesgos}
\begin{columns}[T]
\begin{column}{0.49\textwidth}
\textbf{Etapas} \scriptsize(ordenadas por dependencia, cada una con
entregable verificable)
\smallskip
\scriptsize
\begin{description}
  \item[0. Consolidación] Un solo operador directo, con tests de regresión
    contra los tres solvers y fragmentación regularizada.
  \item[1. No-unicidad] Teorema $J_4$ y espacio nulo; familia de Webster.
    \alert{Parcialmente hecho.}
  \item[2. Identificabilidad] Tabla de todos los parámetros frente a todos los
    conjuntos de datos. ¿Cuánto aporta el tiempo? ¿Y la dimensión radial?
  \item[3. Inversión de Calbuco] Posteriores con datos reales; verificación
    predictiva.
  \item[4. Geometría] Reconstrucción regularizada de $R(z)$; efecto del
    regularizador.
  \item[5. Escritura] Manuscrito $+$ un artículo candidato (capítulos 2 y 3).
\end{description}
\end{column}
\begin{column}{0.49\textwidth}
\textbf{Riesgos y mitigaciones}
\smallskip
\scriptsize
\begin{description}
  \item[Operador caro, MCMC no converge] Emulador (E5); o estadísticos resumen
    con algoritmo de vecindad, como Wong \emph{et al.} (2020).
  \item[\texttt{scikits.odes} no compila] El prototipo sólo usa numpy/scipy;
    portar el estacionario a un DAE puro de scipy.
  \item[Datos insuficientes de Calbuco] Ése \alert{es} el resultado: se reporta
    como análisis de identificabilidad y se complementa con datos sintéticos.
  \item[Los \emph{switches} impiden gradientes] Regularizar la fragmentación
    (E1); si falla, métodos sin derivadas o emulador.
  \item[El alcance crece demasiado] El núcleo son los capítulos 2 y 3, que
    \alert{ya tienen resultados preliminares}; 4 y 5 son incrementales.
\end{description}
\end{column}
\end{columns}
\end{frame}

% ===========================================================================
\section{Preguntas para ustedes}
% ===========================================================================

\begin{frame}{Para el profesor de matemáticas}
\begin{enumerate}
  \item ¿Enfoque \textbf{determinista} (Tikhonov, tasas de convergencia bajo
    condiciones de fuente) o \textbf{bayesiano} (posterior, MCMC, cuantificación
    de incertidumbre)? El capítulo 5 admite ambos; la elección cambia el sabor
    de la tesis.
  \item ¿Vale la pena invertir en el \textbf{método adjunto} (E2), o basta con
    diferencias finitas y un número moderado de parámetros?
  \item ¿Le parece suficiente el resultado de no-unicidad del teorema $J_4$ como
    \textbf{teorema central}, o conviene apuntar a un enunciado de
    \textbf{unicidad condicional} (bajo qué datos adicionales $R(z)$ \emph{sí}
    queda determinada)?
  \item La analogía de \textbf{Webster–Sturm–Liouville} (D2): ¿capítulo propio o
    sección?
\end{enumerate}

\mensaje{Herramientas importables que ya identifiqué: perfiles de verosimilitud
(Raue \emph{et al.}, 2009); teoría espectral inversa (Borg 1946;
Gel'fand–Levitan; Pöschel \& Trubowitz 1987); identificación de coeficientes de
Robin (Chaabane \& Jaoua 1999; Sincich 2007); problemas parabólicos laterales
(Eldén \emph{et al.} 2000); bayesiano en dimensión infinita (Stuart, 2010).}
\end{frame}

% ---------------------------------------------------------------------------
\begin{frame}{Para el profesor de geología}
\begin{enumerate}
  \item Sobre ``recuperar la presión en profundidad'': ¿el objetivo es
    $\Delta P$ en la cámara, o el perfil $P(z)$ completo?
  \item Sobre el contenido de burbujas en profundidad: ¿hay \textbf{serie
    temporal} de porosidad/densidad de piroclastos por unidad estratigráfica
    para Calbuco? Si la hay, \alert{B3 se vuelve el problema inverso más directo
    y mejor condicionado de todo el catálogo}.
  \item ¿Existen mediciones de \textbf{permeabilidad} de los piroclastos?
    Habilitan A4/B4, y la literatura indica que la permeabilidad es el control de
    primer orden.
  \item ¿Qué \textbf{evidencia independiente} hay sobre la geometría del
    conducto de Calbuco (componente lítico de los depósitos, geometría del
    cráter, diques aflorantes)? Serviría como previa y como verificación del
    capítulo 5.
  \item ¿Confirma $T$ y $c_0$? (Ver la advertencia de consistencia.)
\end{enumerate}
\end{frame}

% ---------------------------------------------------------------------------
\begin{frame}{Para ambos}
\begin{block}{La pregunta de alcance}
¿El caso de estudio es \textbf{Calbuco 2015}, o conviene un caso con mejor
cobertura de series de tiempo —por ejemplo un caso efusivo tipo domo, donde la
tasa de extrusión se mide continuamente?
\end{block}

\medskip
El análisis de identificabilidad del experimento 3 sugiere que \alert{los casos
con series de tiempo largas son mucho más informativos}. Eso es un argumento
\emph{técnico}, no de gusto, para la elección del caso.

\medskip
Contraargumento a favor de Calbuco: es chileno, es explosivo (nadie ha invertido
un caso explosivo con modelo físico), tiene petrología publicada de buena
calidad, y ya tenemos el modelo calibrado con sus parámetros.

\vspace{1em}
\begin{exampleblock}{Alternativa}
Usar \textbf{ambos}: un caso efusivo bien muestreado como banco de pruebas
metodológico del capítulo 3, y Calbuco como aplicación del capítulo 4.
\end{exampleblock}
\end{frame}

% ---------------------------------------------------------------------------
\begin{frame}[standout]
Gracias

\vspace{1em}
{\small\normalfont
Todo lo presentado es reproducible:\\
\texttt{tesis/PROPUESTA.md}\quad
\texttt{tesis/ESTADO\_DEL\_ARTE.md}\quad
\texttt{tesis/prototipo/}}
\end{frame}

% ===========================================================================
\appendix
% ===========================================================================

\begin{frame}{Respaldo — el modelo reducido en detalle}
\small
Dominio $z\in[H,z_f]$, con $P(z_f)=P_{frag}$ definida por
$\phi_{eq}(P_{frag})=\phi_{frag}$: es un dato termodinámico, no un parámetro
ajustable. Sobre $z_f$ el flujo es una dispersión gas–piroclastos que el modelo
reducido no resuelve (cortar ahí evita las velocidades de salida no físicas que
da la lubricación en la región fragmentada).

\medskip
\textbf{Cierre termodinámico} (exsolución en equilibrio, isotermo):
$$c_d(P)=\min(c_0,C_1P^\beta),\quad n(P)=\frac{c_0-c_d}{1-c_d},\quad
\rho_g=\frac{P}{R_vT},\quad \frac1\rho=\frac{n}{\rho_g}+\frac{1-n}{\rho_m}.$$
El $\min$ se reemplaza por una versión suavizada para que el operador sea
diferenciable.

\medskip
\textbf{Limitaciones conocidas:} una sola fase con velocidad única (sin
deslizamiento gas–fundido ni escape por permeabilidad); exsolución y
cristalización en equilibrio; isotermo; cámara elástica de un parámetro sin
recarga.

\medskip
Ninguna afecta las conclusiones cualitativas: \alert{la degeneración $J_4$ es
una propiedad estructural de cualquier modelo con fricción tipo Poiseuille.}
\end{frame}

% ---------------------------------------------------------------------------
\begin{frame}{Respaldo — por qué los perfiles de verosimilitud}
\small
La matriz de información de Fisher $F=J^\top\Sigma^{-1}J$ da elipses de
confianza a partir de una \alert{aproximación cuadrática local} del $\chi^2$.

\medskip
Aquí eso falla: el operador directo es fuertemente no lineal en $\Delta P$ (la
sobrepresión compite con la flotabilidad, y su peso relativo cambia con el
propio $\Delta P$). La elipse de Fisher subestima groseramente la incertidumbre.

\medskip
El \textbf{perfil de verosimilitud}
$$\chi^2_{perfil}(\theta_i)=\min_{\theta_j,\,j\neq i}\chi^2(\theta)$$
recorre el valle completo, sin suponer que sea cuadrático. Un perfil que
\alert{no sube} por encima del umbral $\chi^2_{1,0.95}=3.84$ indica un parámetro
\textbf{prácticamente no identificable} con los datos disponibles — que es
exactamente lo que ocurre con $\Delta P$ por abajo.

\medskip
Es el estándar en biología de sistemas desde Raue \emph{et al.} (2009) y se
importa sin cambios.
\end{frame}

% ---------------------------------------------------------------------------
\begin{frame}{Respaldo — la otra no-unicidad (D3)}
\small
Hay \textbf{dos} no-unicidades distintas, y conviene no confundirlas:

\vspace{0.8em}
\begin{columns}[T]
\begin{column}{0.47\textwidth}
\begin{block}{Del problema \emph{directo}}
Para parámetros fijos pueden existir \alert{varios estados estacionarios}
(Melnik \& Sparks 1999: hasta tres tasas de descarga que difieren en órdenes de
magnitud).

\smallskip
Mecanismo: \emph{resistencia diferencial negativa} $dp_{ch}/dq<0$, por
cristalización retardada y por cambio de porosidad debido al escape de gas
(Kozono \& Koyaguchi 2012).

\smallskip
\emph{La literatura sí la conoce.}
\end{block}
\end{column}
\begin{column}{0.47\textwidth}
\begin{block}{Del problema \emph{inverso}}
Para un dato fijo existe una \alert{familia de parámetros y geometrías}
compatibles: $\mathcal{F}$ no es inyectivo.

\smallskip
Mecanismo: el dato ve la geometría sólo a través de un funcional escalar
(teorema $J_4$); el espacio nulo tiene codimensión 1.

\smallskip
\emph{Hasta donde he revisado, no ha sido formalizada en volcanología.}
\end{block}
\end{column}
\end{columns}

\vspace{0.8em}
Ambas empeoran la inversión, y la primera además rompe la unicidad del
\emph{forward}, lo que obliga a un tratamiento cuidadoso (continuación,
selección de rama). Vale la pena tratarlas juntas en el capítulo 2.
\end{frame}

% ---------------------------------------------------------------------------
\begin{frame}[fragile]{Respaldo — estructura de archivos}
\scriptsize
\begin{semiverbatim}
tesis/
  PROPUESTA.md                           \textmd{documento completo}
  ESTADO\_DEL\_ARTE.md                     \textmd{revisión anotada, con DOIs}
  prototipo/
    README.md                            \textmd{cómo reproducir todo}
    modelo\_reducido.py                   \textmd{operador directo con R(z)}
    exp1\_no\_unicidad\_geometria.py
    exp2\_el\_tiempo\_rompe\_la\_degeneracion.py
    exp3\_identificabilidad.py
    exp4\_cuando\_hace\_falta\_el\_transiente.py
    figuras/                             \textmd{figuras de esta presentación}
  presentacion/
    propuesta\_tesis.tex                  \textmd{esta presentación}
    propuesta\_tesis.pdf
\end{semiverbatim}

\vspace{0.8em}
\small
Reproducir todo:
\scriptsize
\begin{semiverbatim}
pip install numpy scipy matplotlib
cd tesis/prototipo
python3 exp1\_no\_unicidad\_geometria.py
python3 exp2\_el\_tiempo\_rompe\_la\_degeneracion.py
python3 exp3\_identificabilidad.py
python3 exp4\_cuando\_hace\_falta\_el\_transiente.py
cd ../presentacion && latexmk -pdf propuesta\_tesis.tex
\end{semiverbatim}
\end{frame}

\end{document}
